\section{本讲目标：从架构动物园中提取默认值}

本节建立全讲方法：不要逐个追随 2024--2025 的模型 tweak，而要从公开 dense model 中提取稳定共识和高风险变量。

先看课程给出的路线图。它没有按模型公司或发布时间组织材料，而是先建立一个现代 Transformer baseline，再把公开模型之间的差异分解为 architecture、hyperparameter 与 stability 三类问题。这样做的好处是：新模型发布时，只需判断它改变了哪一层决策，而不必从头记忆整张配置表。

\begin{figure}[H]
\centering
\includegraphics[width=0.88\textwidth]{slide-001.jpg}
\caption{Slide 2：全讲路线与目标；从动手实现出发，再读取公开模型的共同经验。}
\figsource{CS336 Spring 2026 Lecture 3 官方幻灯片。}
\end{figure}

\begin{knowledgebox}{读图：为什么先讲“如何学习架构”}
这页把 hands-on experience 放在首位，因为架构名词只有落到 tensor shape、参数量、数值范围和运行时上才真正可比较；公开模型报告则是第二层证据，用来发现跨团队重复出现的默认值。后文每个结论都应回答三个问题：它解决什么瓶颈、证据来自质量还是系统指标、换规模或硬件后是否仍成立。
\end{knowledgebox}


\begin{figure}[H]
\centering
\includegraphics[width=0.88\textwidth]{slide-002.jpg}
\caption{Slide 3: Starting point: original Transformer}
\figsource{CS336 Spring 2026 Lecture 3 官方幻灯片。}
\end{figure}

\noindent\textbf{展开说明：}本页用原始 Transformer 作为起点：sinusoidal position embedding、ReLU FFN、post-norm LayerNorm。后续所有现代变体都可以看成对这三个默认值的替换。

与原始版本对照后，Slide 4 把作业中实现的“简单现代版”单独列出。这里真正值得记住的不是四个名词，而是四类迁移：norm 从 residual add 之后移到子层之前；position 从输入端相加改为 Q/K 旋转；FFN 从单支非线性改成 gated branch；线性层去掉 bias。后文会逐项判断这些变化是质量共识、稳定性要求，还是数据移动优化。

\begin{figure}[H]
\centering
\includegraphics[width=0.88\textwidth]{slide-003.jpg}
\caption{Slide 4：课程作业采用的现代 Transformer 变体，以及需要追问的四个设计选择。}
\figsource{CS336 Spring 2026 Lecture 3 官方幻灯片。}
\end{figure}

\begin{importantbox}{baseline 不是 recipe 的终点}
Pre-norm、RoPE、SwiGLU 与 no-bias 是可靠起点，却不能自动决定 FFN ratio、head shape、depth/width、regularization 或推理 attention。真正的 recipe 还必须绑定模型规模、训练数据、上下文长度、并行方式和服务目标；本讲后半段正是在补齐这些条件。
\end{importantbox}

接下来两页用模型发布数量解释为什么不能按“逐模型背配置”的方式学习。第一张图把 2024--2025 年 dense releases 并置，重点看它们虽然有大量微调，但主体仍围绕 LLaMA-like block；第二张图进一步强调发布速度和命名噪声会掩盖共识。读者应把模型名当作证据样本，而不是把每一个新名称都当成全新范式。

\begin{figure}[H]
\centering
\includegraphics[width=0.88\textwidth]{slide-004.jpg}
\caption{Slide 5：2024--2025 年大量 dense model release 与细粒度架构 tweak。}
\figsource{CS336 Spring 2026 Lecture 3 官方幻灯片。}
\end{figure}

\begin{knowledgebox}{读图：模型很多，不等于架构空间同样大}
先比较每个模型是否仍使用 decoder-only Transformer、pre-norm、gated FFN 与 RoPE，再看真正变化的列。常见差异集中在 norm placement、attention variant、position treatment 和少数 ratio；这说明调查的目标应是找“稳定列”和“变化列”。发布数量本身不能证明某个 tweak 有效，只有消融、训练稳定性或运行时证据才能支撑因果判断。
\end{knowledgebox}

\begin{figure}[H]
\centering
\includegraphics[width=0.88\textwidth]{slide-005.jpg}
\caption{Slide 6：模型发布密度使逐篇记忆失效，需要转向经验归纳。}
\figsource{CS336 Spring 2026 Lecture 3 官方幻灯片。}
\end{figure}

\begin{warningbox}{不要用“出现频率”替代实验}
某个设计在很多模型里出现，只能说明它是工程上可接受的 strong prior，不能说明它在你的数据、规模与硬件上最优。频率适合确定实验起点；是否保留仍应由 controlled ablation、稳定性监控和端到端 throughput 决定。
\end{warningbox}


\begin{figure}[H]
\centering
\includegraphics[width=0.88\textwidth]{slide-006.jpg}
\caption{Slide 7: Data-driven architecture study}
\figsource{CS336 Spring 2026 Lecture 3 官方幻灯片。}
\end{figure}

\noindent\textbf{展开说明：}本页说明本讲不是凭偏好讲架构，而是看 dense architectures 的公开数据：哪些共同、哪些变化、哪些能学。

从数据表转向具体章节前，Slide 8 把问题压缩成 coverage map：先看 activation、attention 和 position 这些结构变化，再看 FFN/head/vocab/depth-width 等数值选择，最后处理训练不稳定。这个顺序很重要，因为稳定性技巧往往是在既定架构上控制数值尺度，而不是用来掩盖一个未经验证的 architecture recipe。

\begin{figure}[H]
\centering
\includegraphics[width=0.88\textwidth]{slide-007.jpg}
\caption{Slide 8：本讲覆盖的架构变体、超参数问题与稳定性技巧。}
\figsource{CS336 Spring 2026 Lecture 3 官方幻灯片。}
\end{figure}

\begin{knowledgebox}{阅读路线：先机制，后默认值，再看例外}
Activation、position 和 attention 先回答“计算图如何改变”；FFN ratio、head shape 与 aspect ratio 再回答“规模如何分配”；z-loss、QK norm 等最后回答“数值什么时候失控”。把三层混在一起容易产生错误归因，例如把 RMSNorm 的带宽收益说成纯 FLOPs 收益，或把 GQA 的推理优势误写成训练质量定律。
\end{knowledgebox}


\begin{figure}[H]
\centering
\includegraphics[width=0.88\textwidth]{slide-008.jpg}
\caption{Slide 9: Architecture high level view}
\figsource{CS336 Spring 2026 Lecture 3 官方幻灯片。}
\end{figure}

\noindent\textbf{展开说明：}本页给出大方向：LLaMA-like architecture 占主导，同时出现 QK norm、hybrid attention 等趋势。

\subsection{本章小结}
本节的关键不是记住所有模型名，而是把公开模型的设计选择转化成可迁移判断：共识默认值可以作为起点，例外需要知道原因，涉及系统和推理成本的选择必须结合资源账本讨论。

